\documentclass[11pt,a4paper]{article}

\usepackage{amsmath,amssymb,amsthm,mathtools}
\usepackage[margin=1in]{geometry}
\usepackage{hyperref}
\usepackage{enumitem}

\newtheorem{theorem}{Theorem}[section]
\newtheorem{lemma}[theorem]{Lemma}
\newtheorem{corollary}[theorem]{Corollary}
\newtheorem{proposition}[theorem]{Proposition}
\theoremstyle{definition}
\newtheorem{definition}[theorem]{Definition}
\newtheorem{remark}[theorem]{Remark}
\newtheorem{example}[theorem]{Example}

\newcommand{\R}{\mathbb{R}}
\newcommand{\Z}{\mathbb{Z}}
\newcommand{\N}{\mathbb{N}}
\newcommand{\cJ}{\mathcal{J}}

\title{Extending symplectic embeddings to symplectic automorphisms}
\author{}
\date{}

\begin{document}
\maketitle

\begin{abstract}
Let \((V,\Omega)\) be a separable real Hilbert space equipped with a strong symplectic form
induced by an orthogonal complex structure \(\cJ\). We prove two extension theorems.
First, every bounded symplectic embedding \(T:V\to V\) extends to a symplectic automorphism of
a larger symplectic Hilbert space; the construction is elementary and explicit.
Second, every commuting pair of \emph{isometric} symplectic embeddings extends to a commuting
pair of unitary symplectic automorphisms on a common enlargement. The argument is the
symplectic counterpart of It\^o's conjugation method for commuting isometries: one dilates the
first map by the unilateral defect construction, conjugates the second map along the resulting
unitary group, and dilates once more. No double-commutation hypothesis is used.
\end{abstract}

\tableofcontents

\section{Setup and notation}

Throughout, \(V\) is a real separable Hilbert space with inner product \(g\) and norm
\(\|\cdot\|\). Fix \(\cJ\in\mathcal{B}(V)\) satisfying
\begin{equation}
\label{eq:J}
\cJ^2=-I,\qquad \cJ^*=-\cJ.
\end{equation}
The associated bilinear form
\begin{equation}
\label{eq:Omega}
\Omega(x,y):=g(\cJ x,y)
\end{equation}
is then a \emph{strong} symplectic form: the map \(\flat:V\to V^*\) given by
\(\flat(x)=\Omega(x,\cdot)\) is a topological isomorphism.

The motivating model is the countable orthogonal sum
\[
H=\ell^2(\N;\R^{2d})=\bigoplus_{n\ge 1}\R^{2d},
\qquad
\cJ=\bigoplus_{n\ge 1}J_{2d},
\qquad
\Omega=\bigoplus_{n\ge 1}\omega_0,
\]
where \(\omega_0\) is the standard area form on \(\R^{2d}\) and \(J_{2d}\) is the corresponding
orthogonal complex structure.

\begin{definition}
An operator \(T\in\mathcal{B}(V)\) is a \emph{symplectic embedding} if
\begin{equation}
\label{eq:emb}
\Omega(Tx,Ty)=\Omega(x,y)\qquad\text{for all }x,y\in V,
\end{equation}
equivalently \(T^*\cJ T=\cJ\). It is a \emph{symplectic automorphism} if it is additionally
bijective; in that case \(T^{-1}\in\mathcal{B}(V)\) by the bounded inverse theorem, and we write
\(T\in\mathrm{Sp}(V,\Omega)\).

A symplectic embedding \(T\) is called \emph{isometric} if \(T^*T=I\) (equivalently,
\(\|Tx\|=\|x\|\) for every \(x\in V\)).
\end{definition}

\begin{lemma}
\label{lem:basic}
\begin{enumerate}[label=\textup{(\roman*)}]
\item \(\|\cJ x\|=\|x\|\) for every \(x\in V\).
\item Condition \eqref{eq:emb} holds if and only if \(T^*\cJ T=\cJ\).
\item The form \(\Omega\) is strong.
\end{enumerate}
\end{lemma}

\begin{proof}
From \eqref{eq:J} one has \(\cJ^*\cJ=I\), so
\(\|\cJ x\|^2=g(x,\cJ^*\cJ x)=\|x\|^2\). Expanding the definitions gives
\(\Omega(Tx,Ty)=g(T^*\cJ Tx,y)\) and \(\Omega(x,y)=g(\cJ x,y)\), which proves~(ii).
For~(iii), note that \(\flat=\iota_g\circ\cJ\) where \(\iota_g:V\to V^*\) is the Riesz map.
\end{proof}

\section{Left inverse, closed range, and the symplectic complement}

\begin{lemma}
\label{lem:left}
Let \(T\) be a symplectic embedding. Then:
\begin{enumerate}[label=\textup{(\roman*)}]
\item \(\|Tx\|\ge\|x\|/\|T\|\) for all \(x\in V\); in particular \(T\) is injective;
\item \(R:=\operatorname{ran}T\) is closed;
\item the operator \(L:=-\cJ T^*\cJ\) satisfies \(LT=I\) and \(\|L\|\le\|T\|\);
\item \(E:=TL\) is a bounded projection onto \(R\), and \(F:=I-E\) satisfies \(F|_R=0\).
\end{enumerate}
\end{lemma}

\begin{proof}
Using Lemma~\ref{lem:basic},
\[
\|x\|=\|\cJ x\|=\|T^*\cJ Tx\|\le\|T\|\,\|Tx\|,
\]
which is~(i). If \(Tx_n\to z\), then \((Tx_n)\) is Cauchy, hence so is \((x_n)\) by~(i); if
\(x_n\to x\), continuity yields \(Tx=z\), so \(R\) is closed. Next,
\[
LT=(-\cJ T^*\cJ)T=-\cJ(T^*\cJ T)=-\cJ^2=I,
\]
and \(\|L\|\le\|T\|\) because \(\|\cJ\|=1\). Finally \(E^2=T(LT)L=TL=E\) and
\(\operatorname{ran}E=R\), which proves~(iv).
\end{proof}

The operator \(L\) is the adjoint of \(T\) with respect to \(\Omega\). Indeed,
\[
\Omega(Lx,y)=g(\cJ Lx,y)=g(\cJ(-\cJ T^*\cJ)x,y)=g(T^*\cJ x,y)=g(\cJ x,Ty)=\Omega(x,Ty).
\]
Consequently
\[
N:=\ker L=\{v\in V:\Omega(v,Tx)=0\text{ for all }x\in V\}=R^{\perp_\Omega}.
\]

\begin{lemma}
\label{lem:N}
Let \(T\) be a symplectic embedding, \(R=\operatorname{ran}T\), and \(N=R^{\perp_\Omega}\). Then:
\begin{enumerate}[label=\textup{(\roman*)}]
\item \(N=\cJ^{-1}(R^{\perp_g})\) is closed, and \(V=R\oplus N\) as a topological direct sum with
\(\Omega(R,N)=0\);
\item both \((R,\Omega|_R)\) and \((N,\Omega|_N)\) are strong symplectic Hilbert spaces;
\item \(T:V\to R\) is a symplectic isomorphism;
\item \(T\) is surjective if and only if \(N=\{0\}\).
\end{enumerate}
\end{lemma}

\begin{proof}
The identity \(N=\cJ^{-1}(R^{\perp_g})\) is immediate from
\(\Omega(v,r)=g(\cJ v,r)\). Let \(P_R\) be the orthogonal projection onto \(R\) and set
\(B:=P_R\cJ|_R:R\to R\). Then \(g(Br,s)=\Omega(r,s)\) for \(r,s\in R\). We temporarily assume
the nondegeneracy assertion in~(ii) for \(R\) (proved independently in the next paragraph via
the isomorphism \(T:V\to R\)), so \(B\) is bijective. The operator
\(\widetilde{E}:=B^{-1}P_R\cJ\) is then a bounded idempotent with range \(R\) and kernel
\(N\), whence \(V=R\oplus N\) topologically. Comparing kernels and ranges with
Lemma~\ref{lem:left}(iv) shows \(\widetilde{E}=E=TL\). By definition of \(N\) one has
\(\Omega(R,N)=0\).

To see that \(\Omega|_R\) is strong, observe that \(T:V\to R\) is a linear bijection
preserving \(\Omega\) (Lemma~\ref{lem:left} and \eqref{eq:emb}), and \(\Omega\) is strong on
\(V\), so the transported form on \(R\) is strong. In particular \(B\) is an isomorphism, as
used above.

Now suppose \(n\in N\) and \(\Omega(n,N)=0\). Then also \(\Omega(n,R)=0\), so
\(\Omega(n,V)=0\), hence \(n=0\). Thus \(\Omega|_N\) is nondegenerate. Given \(\phi\in N^*\),
the functional \(\tilde\phi:=\phi\circ F\) lies in \(V^*\). Choose \(x\in V\) with
\(\flat(x)=\tilde\phi\), and write \(x=r+n\) according to \(V=R\oplus N\). For every
\(m\in N\),
\[
\phi(m)=\Omega(r+n,m)=\Omega(n,m),
\]
so \(\flat_N(n)=\phi\). Therefore \(\flat_N\) is bijective, and the open mapping theorem
makes it a topological isomorphism.

Assertion~(iii) is Lemma~\ref{lem:left} together with \eqref{eq:emb}, and~(iv) is the
equivalence \(R=V\Longleftrightarrow N=\{0\}\).
\end{proof}

\begin{lemma}
\label{lem:isometric-orthogonal}
If \(T\) is an isometric symplectic embedding, then \(N=R^{\perp_g}\). In particular the
decomposition \(V=R\oplus N\) is orthogonal, and \(\cJ\) leaves both \(R\) and \(N\) invariant.
\end{lemma}

\begin{proof}
From \(T^*T=I\) and \(T^*\cJ T=\cJ\) one obtains \(T^*=-\cJ T^*\cJ=L\). Taking adjoints and
using \(\cJ^*=-\cJ\) yields \(T=-\cJ T\cJ\), or equivalently \(\cJ T=-T\cJ\). Hence
\[
\cJ(R)=\cJ T(V)=T\cJ(V)=R,
\]
and therefore \(\cJ(R^{\perp_g})=R^{\perp_g}\) as well. It follows that
\[
R^{\perp_\Omega}=\cJ^{-1}(R^{\perp_g})=R^{\perp_g}.
\]
\end{proof}

\section{Extension of a single symplectic embedding}

\begin{theorem}
\label{thm:A}
Let \(T:V\to V\) be a symplectic embedding. Then there exist a symplectic Hilbert space
\((W,\widehat{\Omega})\) containing \((V,\Omega)\) as a closed symplectic subspace and a
symplectic automorphism \(\widetilde{T}:W\to W\) such that \(\widetilde{T}|_V=T\).
\end{theorem}

\begin{proof}
If \(N=\{0\}\), then \(T\) is already a symplectic automorphism by Lemma~\ref{lem:N}(iv), and
there is nothing to prove. Assume henceforth that \(N\neq\{0\}\).

For each integer \(k\ge 1\) let \((N_k,g_k,\Omega_k)\) be a copy of \((N,g|_N,\Omega|_N)\), and
fix a linear isomorphism \(\iota_k:N\to N_k\) that is unitary for the Hilbert structures and
satisfies \(\Omega_k(\iota_k n,\iota_k n')=\Omega(n,n')\) for all \(n,n'\in N\). Define
\begin{equation}
\label{eq:W}
W:=V\,\oplus\,\bigoplus_{k\ge 1}N_k
\end{equation}
as an orthogonal direct sum of Hilbert spaces: its elements are sequences
\((v;n_1,n_2,\dots)\) with \(\|v\|^2+\sum_{k\ge 1}\|n_k\|^2<\infty\). Equip \(W\) with the
continuous bilinear form
\begin{equation}
\label{eq:Omegahat}
\widehat{\Omega}\bigl((v;n_\bullet),(v';n'_\bullet)\bigr)
:=\Omega(v,v')+\sum_{k\ge 1}\Omega(n_k,n'_k).
\end{equation}
Continuity of \(\widehat{\Omega}\) follows from \(|\Omega(a,b)|\le\|a\|\|b\|\) and the
Cauchy--Schwarz inequality on \(\ell^2\). The associated map
\(\widehat{\flat}:W\to W^*\) is the orthogonal direct sum
\(\flat\oplus\bigoplus_k\flat_N\). Each summand is a topological isomorphism
(Lemma~\ref{lem:N}(ii)), and the operator norms of the inverses are uniformly bounded (each
copy of \(N\) carries the same symplectic structure). Hence \(\widehat{\flat}\) is a
topological isomorphism, so \(\widehat{\Omega}\) is strong. The isometric embedding
\(v\mapsto(v;0,0,\dots)\) realizes \(V\) as a closed symplectic subspace of \(W\).

Define \(\widetilde{T}:W\to W\) by
\begin{equation}
\label{eq:Ttil}
\widetilde{T}(v;n_1,n_2,n_3,\dots)
=\bigl(Tv+\iota_1^{-1}n_1;\;
\iota_1\iota_2^{-1}n_2,\;
\iota_2\iota_3^{-1}n_3,\;\dots\bigr).
\end{equation}
Each factor \(\iota_k\iota_{k+1}^{-1}\) is unitary, and \(\iota_1^{-1}\) is unitary, so
\begin{align*}
\|\widetilde{T}u\|^2
&=\|Tv+\iota_1^{-1}n_1\|^2+\sum_{k\ge 1}\|n_{k+1}\|^2\\
&\le 2\|T\|^2\|v\|^2+2\|n_1\|^2+\sum_{k\ge 1}\|n_{k+1}\|^2
\le\bigl(2\|T\|^2+2\bigr)\|u\|^2.
\end{align*}
Thus \(\widetilde{T}\in\mathcal{B}(W)\). If \(n_\bullet=0\), then
\(\widetilde{T}(v;0)=(Tv;0,\dots)\), so \(\widetilde{T}|_V=T\).

We next verify that \(\widetilde{T}\) preserves \(\widehat{\Omega}\). Let
\(u=(v;n_\bullet)\) and \(u'=(v';n'_\bullet)\). Then
\begin{align*}
\widehat{\Omega}(\widetilde{T}u,\widetilde{T}u')
&=\Omega\bigl(Tv+\iota_1^{-1}n_1,\,Tv'+\iota_1^{-1}n'_1\bigr)
+\sum_{k\ge 1}\Omega(n_{k+1},n'_{k+1}).
\end{align*}
Since \(Tv\in R\), \(\iota_1^{-1}n_1\in N\), and \(\Omega(R,N)=0\), the cross terms vanish.
Symplecticity of \(T\) and of \(\iota_1\) therefore yields
\[
\Omega(Tv,Tv')=\Omega(v,v'),\qquad
\Omega(\iota_1^{-1}n_1,\iota_1^{-1}n'_1)=\Omega(n_1,n'_1).
\]
Adding the shifted tail sum produces
\(\widehat{\Omega}(\widetilde{T}u,\widetilde{T}u')=\widehat{\Omega}(u,u')\).

Injectivity of \(\widetilde{T}\) follows at once: if \(\widetilde{T}u=0\), then for every
\(u'\in W\)
\[
\widehat{\Omega}(u,u')=\widehat{\Omega}(\widetilde{T}u,\widetilde{T}u')=0,
\]
so \(u=0\) by strongness of \(\widehat{\Omega}\).

For surjectivity, take an arbitrary \(w=(z;m_1,m_2,\dots)\in W\). Set
\(n_{k+1}:=\iota_{k+1}\iota_k^{-1}m_k\) for \(k\ge 1\); then
\(\sum\|n_{k+1}\|^2=\sum\|m_k\|^2<\infty\). Decompose \(z=z_R+z_N\) according to
\(V=R\oplus N\), and put \(v:=T^{-1}z_R\in V\) together with \(n_1:=\iota_1 z_N\). The vector
\(u=(v;n_1,n_2,\dots)\) belongs to \(W\) and satisfies \(\widetilde{T}u=w\).

Thus \(\widetilde{T}\) is a bounded bijective endomorphism of the Hilbert space \(W\), so
\(\widetilde{T}^{-1}\) is bounded by the bounded inverse theorem. The inverse preserves
\(\widehat{\Omega}\) because \(\widetilde{T}\) does. Therefore
\(\widetilde{T}\in\mathrm{Sp}(W,\widehat{\Omega})\).
\end{proof}

\begin{remark}
\label{rem:inverse}
Writing \(E=TL\) and \(F=I-E\), the inverse is given explicitly by
\begin{equation}
\label{eq:Tinv}
\widetilde{T}^{-1}(z;m_1,m_2,\dots)
=\bigl(Lz;\;
\iota_1 Fz,\;
\iota_2\iota_1^{-1}m_1,\;
\iota_3\iota_2^{-1}m_2,\;\dots\bigr).
\end{equation}
Indeed, a direct substitution into \eqref{eq:Ttil} recovers the identity on \(W\), using
\(TLz+F z=z\) and \(\iota_k\iota_{k+1}^{-1}\iota_{k+1}\iota_k^{-1}=\mathrm{id}\).
Iterating on vectors of \(V\) yields, for every \(k\ge 1\) and \(v\in V\),
\begin{equation}
\label{eq:neg}
\widetilde{T}^{-k}v
=\bigl(L^kv;\;
\iota_1 FL^{k-1}v,\;
\dots,\;
\iota_k Fv,\;
0,0,\dots\bigr).
\end{equation}
In particular, if \(n\in N\) then \(Ln=0\) and \(Fn=n\), so
\begin{equation}
\label{eq:negN}
\widetilde{T}^{-k}n=(0;\,0,\dots,0,\iota_k n,0,\dots).
\end{equation}
\end{remark}

\begin{remark}
\label{rem:unitary}
If \(T\) is an isometric symplectic embedding, then the operator \(\widetilde{T}\) of
Theorem~\ref{thm:A} is unitary. Indeed, Lemma~\ref{lem:isometric-orthogonal} provides an
orthogonal decomposition \(V=R\oplus N\) with \(\|Tv\|=\|v\|\) and with
\(\iota_1^{-1}n_1\perp Tv\). Consequently
\begin{align*}
\|\widetilde{T}u\|^2
&=\|Tv\|^2+\|\iota_1^{-1}n_1\|^2+\sum_{k\ge 1}\|n_{k+1}\|^2
=\|v\|^2+\sum_{k\ge 1}\|n_k\|^2=\|u\|^2.
\end{align*}
Combined with the surjectivity already proved, \(\widetilde{T}\) is unitary. Moreover, the
same computation that gave \(\cJ T=-T\cJ\) extends componentwise to show that
\(\widetilde{T}\) anticommutes with the block-diagonal extension of \(\cJ\) to \(W\), so
\(\widetilde{T}\) is orthosymplectic.
\end{remark}

\begin{example}
If \(T\) is the unilateral block shift on \(\bigoplus_{n\ge 1}\R^{2d}\), then \(N\) is the
first summand and \(\widetilde{T}\) is unitarily equivalent to the bilateral block shift on
\(\bigoplus_{n\in\Z}\R^{2d}\).
\end{example}

\section{Extension of a commuting isometric pair}

Throughout this section, \(T\) and \(T'\) denote \emph{isometric} symplectic embeddings of
\(V\) into itself that commute: \(TT'=T'T\). We do \emph{not} assume that they doubly
commute (i.e., we do not assume \(T^*T'=T'T^*\)). Commuting isometries need not doubly
commute in general, and the argument below nowhere requires this stronger relation.

\begin{theorem}
\label{thm:B}
Let \(T,T':V\to V\) be commuting isometric symplectic embeddings. Then there exist a
symplectic Hilbert space \((W,\widehat{\Omega})\) containing \((V,\Omega)\) as a closed
symplectic subspace and commuting unitary symplectic automorphisms
\(\widetilde{T},\widetilde{T}'\) of \(W\) such that \(\widetilde{T}|_V=T\) and
\(\widetilde{T}'|_V=T'\).
\end{theorem}

\begin{proof}
By Theorem~\ref{thm:A} and Remark~\ref{rem:unitary} there exist a symplectic Hilbert space
\((W_1,\widehat{\Omega}_1)\) and a unitary symplectic automorphism \(U\) of \(W_1\) with
\(U|_V=T\). Replacing \(W_1\) by the closed subspace
\begin{equation}
\label{eq:minimal}
W_1=\overline{\operatorname{span}}\{\,U^k v:k\in\Z,\ v\in V\,\}
\end{equation}
if necessary, we may assume that the extension is minimal; the subspace
\eqref{eq:minimal} reduces \(U\) and inherits a strong symplectic form from
\(\widehat{\Omega}_1\).

\medskip
\noindent\textbf{Construction of a densely defined intertwiner.}
Let \(\mathcal{D}_0:=\operatorname{span}\{U^k v:k\in\Z,\ v\in V\}\), which is dense in
\(W_1\) by minimality. Define a linear map \(W_0:\mathcal{D}_0\to W_1\) on generators by
\begin{equation}
\label{eq:W0}
W_0(U^k v):=U^k(T'v),\qquad k\in\Z,\ v\in V,
\end{equation}
and extend by linearity. To see that \(W_0\) is well-defined and isometric, it suffices to
prove the inner-product identity
\begin{equation}
\label{eq:inner}
g\bigl(U^k(T'x),\,U^\ell(T'y)\bigr)=g\bigl(U^k x,\,U^\ell y\bigr)
\end{equation}
for all \(k,\ell\in\Z\) and all \(x,y\in V\). Indeed, if a linear combination
\(\xi=\sum_i U^{k_i}x_i\) vanishes, then
\[
\biggl\|W_0\xi\biggr\|^2
=\sum_{i,j}g\bigl(U^{k_i}(T'x_i),\,U^{k_j}(T'x_j)\bigr)
=\sum_{i,j}g\bigl(U^{k_i}x_i,\,U^{k_j}x_j\bigr)=\|\xi\|^2=0,
\]
so \(W_0\xi=0\); thus \(W_0\) is independent of the representation of its argument. The same
computation shows \(\|W_0\xi\|=\|\xi\|\) on \(\mathcal{D}_0\).

Since \(U\) is unitary, both sides of \eqref{eq:inner} depend only on the difference
\(k-\ell\). Replacing \((k,\ell)\) by \((k-\ell,0)\) when \(k\ge\ell\), and using unitarity to
pass from the case \(k<\ell\) to the case \(k>\ell\) upon swapping the two arguments, one
reduces \eqref{eq:inner} to the family of identities
\begin{align}
g\bigl(U^m(T'x),\,T'y\bigr)&=g\bigl(U^m x,\,y\bigr),
\label{eq:pos}\\
g\bigl(T'x,\,U^m(T'y)\bigr)&=g\bigl(x,\,U^m y\bigr),
\label{eq:neg2}
\end{align}
for all integers \(m\ge 0\) and all \(x,y\in V\). Identity \eqref{eq:neg2} is the same as
\eqref{eq:pos} with the roles of the two arguments interchanged, so it is enough to prove
\eqref{eq:pos}.

For \(m=0\), \eqref{eq:pos} is the isometry relation \(g(T'x,T'y)=g(x,y)\). Now fix
\(m\ge 1\). On the original space \(V\) one has \(TT'=T'T\) and \(U|_V=T\), whence by a
trivial induction
\begin{equation}
\label{eq:poscomm}
U^m(T'v)=T'(U^m v)\qquad\text{for all }v\in V\text{ and all }m\ge 0.
\end{equation}
(Negative powers of \(U\) are not needed on \(V\) at this stage.) Consequently
\begin{align*}
g\bigl(U^m(T'x),\,T'y\bigr)
&=g\bigl(T'(U^m x),\,T'y\bigr)
=g\bigl(U^m x,\,y\bigr),
\end{align*}
where the second equality uses that \(T'\) is isometric. This is \eqref{eq:pos}.

Thus \(W_0\) is a well-defined isometric linear map on the dense subspace \(\mathcal{D}_0\).
It extends uniquely by uniform continuity to an isometry \(W\in\mathcal{B}(W_1)\) of
\(W_1\) into itself.

\medskip
\noindent\textbf{Extension and intertwining properties.}
For every \(v\in V\) one has \(Wv=W_0v=T'v\), so \(W|_V=T'\). Moreover, if \(U^kv\in\mathcal{D}_0\),
then
\[
WU(U^kv)=W(U^{k+1}v)=U^{k+1}(T'v)=U\bigl(U^k(T'v)\bigr)=UW(U^kv).
\]
By denseness of \(\mathcal{D}_0\) and continuity, \(WU=UW\).

\medskip
\noindent\textbf{Preservation of the symplectic form.}
Write \(\widehat{\Omega}\) for \(\widehat{\Omega}_1\). Let \(k,\ell\in\Z\) and \(x,y\in V\).
Then
\begin{align*}
\widehat{\Omega}\bigl(W U^k x,\,W U^\ell y\bigr)
&=\widehat{\Omega}\bigl(U^k T'x,\,U^\ell T'y\bigr).
\end{align*}
Since \(U\) preserves \(\widehat{\Omega}\),
\[
\widehat{\Omega}\bigl(U^k T'x,\,U^\ell T'y\bigr)
=\widehat{\Omega}\bigl(T'x,\,U^{\ell-k}T'y\bigr).
\]
Set \(m:=\ell-k\). If \(m\ge 0\), then \(U^m\) agrees with \(T^m\) on \(V\), and
\eqref{eq:poscomm} gives \(U^m T'y=T'T^my\). Therefore
\begin{align*}
\widehat{\Omega}\bigl(T'x,\,U^m T'y\bigr)
&=\Omega(T'x,\,T'T^my)
=\Omega(x,\,T^my)
=\widehat{\Omega}\bigl(x,\,U^m y\bigr)
=\widehat{\Omega}\bigl(U^k x,\,U^\ell y\bigr),
\end{align*}
where the second equality uses that \(T'\) preserves \(\Omega\). If \(m=-p\) with \(p>0\),
unitarity of \(U\) moves the negative power to the first argument:
\begin{align*}
\widehat{\Omega}\bigl(T'x,\,U^{-p}T'y\bigr)
&=\widehat{\Omega}\bigl(U^p T'x,\,T'y\bigr)
=\widehat{\Omega}\bigl(T'T^px,\,T'y\bigr)
=\Omega(T^px,\,y)\\
&=\widehat{\Omega}\bigl(U^p x,\,y\bigr)
=\widehat{\Omega}\bigl(x,\,U^{-p}y\bigr)
=\widehat{\Omega}\bigl(U^k x,\,U^\ell y\bigr).
\end{align*}
Thus \(\widehat{\Omega}(W\xi,W\eta)=\widehat{\Omega}(\xi,\eta)\) for all
\(\xi,\eta\in\mathcal{D}_0\). By continuity, \(W\) preserves \(\widehat{\Omega}\) on all of
\(W_1\). In other words, \(W\) is an isometric symplectic embedding of \(W_1\) into itself,
and \(W\) commutes with the unitary symplectic automorphism \(U\).

\medskip
\noindent\textbf{Unitary dilation of \(W\).}
Apply Theorem~\ref{thm:A} in the isometric situation of Remark~\ref{rem:unitary} to the
operator \(W\) on \(W_1\). One obtains a symplectic Hilbert space \((W,\widehat{\Omega})\)
containing \(W_1\) and a unitary symplectic automorphism \(\widetilde{T}'\) of \(W\) with
\(\widetilde{T}'|_{W_1}=W\). Replacing \(W\) by the closed span of
\(\{(\widetilde{T}')^k w_1:k\in\Z,\ w_1\in W_1\}\) if needed, we may assume that this second
dilation is likewise minimal.

\medskip
\noindent\textbf{Extension of \(U\) by conjugation.}
Let \(\mathcal{D}_1:=\operatorname{span}\{(\widetilde{T}')^k w_1:k\in\Z,\ w_1\in W_1\}\),
dense in \(W\) by the minimality just arranged. Define \(U_0:\mathcal{D}_1\to W\) on
generators by
\begin{equation}
\label{eq:U0}
U_0\bigl((\widetilde{T}')^k w_1\bigr):=(\widetilde{T}')^k(U w_1),
\qquad k\in\Z,\ w_1\in W_1,
\end{equation}
and extend linearly. The same reasoning as for \(W_0\) shows that \(U_0\) is well-defined and
isometric, once one verifies
\begin{equation}
\label{eq:inner2}
g\bigl((\widetilde{T}')^k(U\xi),\,(\widetilde{T}')^\ell(U\eta)\bigr)
=g\bigl((\widetilde{T}')^k\xi,\,(\widetilde{T}')^\ell\eta\bigr)
\end{equation}
for \(k,\ell\in\Z\) and \(\xi,\eta\in W_1\). Because \(\widetilde{T}'\) is unitary, it is
enough to treat \(\ell=0\) and \(k=m\ge 0\), together with the symmetric case obtained by
swapping arguments. For \(m=0\) the claim is that \(U\) is isometric on \(W_1\), which is
true. For \(m\ge 1\) we use that \(\widetilde{T}'|_{W_1}=W\) commutes with \(U\), whence by
induction
\[
(\widetilde{T}')^m(U\xi)=U\bigl((\widetilde{T}')^m\xi\bigr)
\qquad\text{for all }\xi\in W_1\text{ and }m\ge 0.
\]
Therefore
\begin{align*}
g\bigl((\widetilde{T}')^m(U\xi),\,U\eta\bigr)
&=g\bigl(U((\widetilde{T}')^m\xi),\,U\eta\bigr)
=g\bigl((\widetilde{T}')^m\xi,\,\eta\bigr),
\end{align*}
as required. Extending by continuity, one obtains an isometry \(\widetilde{U}\in\mathcal{B}(W)\)
satisfying \(\widetilde{U}|_{W_1}=U\) and \(\widetilde{U}\,\widetilde{T}'=\widetilde{T}'\,\widetilde{U}\).

The identical computation as in the form-preservation paragraph above, with the pair
\((U,W)\) replaced by the pair \((\widetilde{U},\widetilde{T}')\) and with \(W_1\) playing the
role formerly played by \(V\), shows that \(\widetilde{U}\) preserves \(\widehat{\Omega}\).
(Explicitly: for \(m\ge 0\),
\[
\widehat{\Omega}\bigl((\widetilde{T}')^m U\xi,\,U\eta\bigr)
=\widehat{\Omega}\bigl(U(\widetilde{T}')^m\xi,\,U\eta\bigr)
=\widehat{\Omega}\bigl((\widetilde{T}')^m\xi,\,\eta\bigr),
\]
while for negative exponents one again moves powers with unitarity of \(\widetilde{T}'\).)

\medskip
\noindent\textbf{Surjectivity of \(\widetilde{U}\).}
The range of the isometry \(\widetilde{U}\) is a closed subspace of \(W\). It contains
\(W_1\) because \(\widetilde{U}|_{W_1}=U\) is unitary on \(W_1\), hence surjective onto
\(W_1\). Moreover \(\operatorname{ran}\widetilde{U}\) is invariant under \(\widetilde{T}'\):
if \(w=\widetilde{U}z\), then
\[
\widetilde{T}'w=\widetilde{T}'\widetilde{U}z=\widetilde{U}(\widetilde{T}'z)\in\operatorname{ran}\widetilde{U}.
\]
Since \(\widetilde{T}'\) is unitary, \(\operatorname{ran}\widetilde{U}\) is likewise invariant
under \((\widetilde{T}')^{-1}\). Therefore \(\operatorname{ran}\widetilde{U}\) contains every
vector \((\widetilde{T}')^k w_1\) with \(k\in\Z\) and \(w_1\in W_1\). By minimality of the
dilation of \(W\), one concludes that \(\operatorname{ran}\widetilde{U}=W\). An isometric
surjective operator on a Hilbert space is unitary, so \(\widetilde{U}\) is unitary.

Set \(\widetilde{T}:=\widetilde{U}\) and retain the notation \(\widetilde{T}'\) for the
unitary already constructed. Then \(\widetilde{T}\) and \(\widetilde{T}'\) are commuting
unitary symplectic automorphisms of \(W\) extending \(T\) and \(T'\) respectively.
\end{proof}

\begin{corollary}
\label{cor:family}
The same conclusion holds for any finite or countably infinite family of mutually commuting
isometric symplectic embeddings. One proceeds by induction on the number of operators (or by
exhaustion in the countable case): at stage \(n\), the first \(n\) operators have been extended
to commuting unitary symplectic automorphisms of a common space \(W^{(n)}\), and the next
isometry is conjugated onto the minimal unitary dilation of one of them (or of a fixed
amalgamation) exactly as in the proof of Theorem~\ref{thm:B}.
\end{corollary}

\begin{example}
If \(T\) is the unilateral block shift on \(\bigoplus_{n\ge 1}\R^{2d}\) and \(T'=T^m\) for
some integer \(m\ge 1\), then both maps are isometric symplectic embeddings, they commute, and
Theorem~\ref{thm:B} recovers the bilateral block shift together with its \(m\)-th power on
\(\bigoplus_{n\in\Z}\R^{2d}\).
\end{example}

\section{Remarks on the non-isometric commuting case}

Theorem~\ref{thm:A} requires no isometry assumption: an arbitrary bounded symplectic embedding
extends to a symplectic automorphism. The conjugation method of Theorem~\ref{thm:B}, however,
uses isometry of \(T'\) (and unitarity of the dilation of \(T\)) in an essential way. Two
obstructions appear if one drops isometry.

First, if \(T\) and \(T'\) are commuting symplectic embeddings that are not isometric, the
defect space \(N=(\operatorname{ran}T)^{\perp_\Omega}\) need not be invariant under \(T'\).
A concrete illustration on \(\bigoplus_{n\ge 1}\R^{2}\) is given by the commuting pair
\begin{align*}
T(x_1,x_2,x_3,\dots)&=(0,\,2x_1,\,\tfrac12 x_2,\,2x_3,\,\tfrac12 x_4,\,\dots),\\
T'(x_1,x_2,x_3,\dots)&=(0,\,3x_1,\,\tfrac13 x_2,\,3x_3,\,\tfrac13 x_4,\,\dots),
\end{align*}
each of which is a (non-surjective, non-isometric) symplectic embedding for the orthogonal sum
of area forms. Here \(N\) is the first summand, while \(T'N\) lies in the second summand and
meets \(N\) only at \(0\). Consequently one cannot define a fibrewise extension of \(T'\) on
the defect slots of the construction in Theorem~\ref{thm:A}.

Second, if one nevertheless forms the minimal automorphism extension \(U=\widetilde{T}\) of
\(T\) furnished by Theorem~\ref{thm:A} and attempts to define a densely defined intertwiner by
the same formula \(U^k v\mapsto U^k(T'v)\), the inner-product identities that produced an
isometric (hence bounded) map on the dense subspace fail: one no longer has
\(\|U^k(T'x)\|=\|U^k x\|\). Boundedness of the conjugated operator on the whole minimal
dilation may therefore fail.

A joint commuting extension to symplectic automorphisms is nevertheless expected on general
grounds after complexification, by viewing \(T\) and \(T'\) as commuting \(J\)-norm-preserving
operators on a Kre\u{\i}n space and invoking the corresponding dilation theory
\cite{KreinDil}. An alternative elementary approach would be to pass to a carefully chosen
non-minimal enlargement on which a bounded intertwiner can be arranged; we do not develop that
construction here. The isometric case treated in Theorem~\ref{thm:B} is the natural symplectic
analogue of It\^o's theorem and is the setting in which the conjugation method closes in a
fully elementary way.

\begin{thebibliography}{5}

\bibitem{Ando}
T.~And\^o,
\textit{On a pair of commutative contractions},
Acta Sci.\ Math.\ (Szeged) \textbf{24} (1963), 88--90.

\bibitem{Ito}
T.~It\^o,
\textit{On the commutative family of subnormal operators},
J.\ Fac.\ Sci.\ Hokkaido Univ.\ Ser.\ I \textbf{12} (1958), 109--127.

\bibitem{KreinDil}
F.~H.~Szafraniec et al.,
\textit{Regular dilations on Kre\u{\i}n spaces},
Canad.\ J.\ Math.\ (2025).

\bibitem{Parthasarathy}
K.~R.~Parthasarathy,
\textit{Symplectic dilations, Gaussian states and Gaussian channels},
arXiv:1405.6476.

\bibitem{SzNagy}
B.~Sz.-Nagy and C.~Foia\c{s},
\textit{Harmonic analysis of operators on Hilbert space},
North-Holland/Akad\'emiai Kiad\'o, 1970.

\end{thebibliography}

\end{document}
